\scanpage{136}
% archive foliation: 135
La codim coh de $M$ sur $B$ est celle de $\bar M$ sur $\bar B$. On en déduit: sous les conditions de cor.2, la codim. coh. de $M$ sur $B$ est identique à celle de $\bar M$ sur $\bar A$ (l'énoncer en corollaire; je m'aperçois que le cor.3 précédent devrait s'énoncer ainsi, avec : par conséquent $\dim_B M\quad\dim A$). On dit qu'un module de type fini sur un anneau local noethérien est Coh Mac si sa dim est égale à sa codim coh d'où un certain nombre de conséquences du lemme (qu'on peut donner en corollaires): Coh. Mac. est invariant par complétion, par changement fini d'anneau, enfin sous les conditions du cor.2 précédent \deleted{\unsure{xx}}, $M$ est Coh Mac sur $B$ si et seulement si $\bar M$ l'est sur $\bar A$. Enfin rappeler que si l'anneau local est régulier, Cohen Mac sur lui signifie libre. En fait, il faudrait sans doute un lemme spécial pour les Coh. Mac, en tous points parallèle (avec ses corollaires) au lemme sur la cocodim, mais contenant en plus le coup du régulier, et peut-être (pour avoir tout ensemble) l'équidimensionnalité. Pour l'instant, il n'était guère question de platitude dans ce numéro, mais:

\uline{Proposition} Soient $A\quad B$, \deleted{\unsure{M comme}\illegible{3}} un hom local d'anneaux locaux noethériens, $M$ un $B$-module de type fini. Pour que $M$ soit $A$-plat, il faut et il suffit que $\bar M$ (le complété du $B$-module $M$ pour sa topologie naturelle) soit plat sur $\bar A$.

Il faut évidemment noter que $\bar M$ (qui n'est pas en général le complété $\uline{m}$-adique de $A$, puisqu'onne fait pas d'hypothèse quasi-finie) est cependant un module sur $\bar A$. -- Démonstration: On applique le critère de prop.13, en notant que $\operatorname{Tor}_1^A(\uline{k},M)$ est un $B$-module de type fini, et que son complété est $\operatorname{Tor}_1^{\bar A}(\uline{k},\bar M)$.

\uline{Corollaire 1} Supposons de plus $M$ qu f sur $A$. Alors $M$ plat sur $A$ équivaut à $\bar M$ libre sur $\bar A$.

\uline{Corollaire 2} Supposons $M$ comme dans le cor.1 et $A$-plat, $\ne0$. Alors
\[
\dim_B M=\dim A\qquad \operatorname{cocodim}_B M=\operatorname{cocodim}_{\deleted{\unsure{\bar A}}}\bar A
\]
En effet, dim et cocodim sont invariants par complétion, d'autre part par changement raisonnable d'anneau, ce qui nous ramène d'abord au cas où $B$ lui-même\unsure{x} est qu f sur $A$, puis par complétion au cas où $A$ et $B$ complets (donc $B$ fini sur $A$) puis encore au\deleted{\unsure{cas}} $A=B$, enfin on note que la dim -- cocodim d'une somme directe est le sup resp. inf des dim \unsure{re}

